\input{estilo}
% La pizarra: dos productores encolan; el consumidor de arriba tomó m₁ y
% falló (la cruz roja); los otros dos tiran de la cola; de m₁ baja el visibility
% timer de 60 segundos. Al costado, "efecto balanceo de carga".
\begin{document}
\begin{figura}\begin{tikzpicture}[
    q/.style={caja, minimum width=11mm, minimum height=11mm},
    m/.style={draw=tinta, fill=white, minimum width=4mm, minimum height=5.5mm, inner sep=0pt},
  ]
  \node[cola, minimum height=40mm, minimum width=11mm] (q) at (3.4,0) {};
  \node[m] (m) at (3.4,0) {};
  \node[etiq, anchor=east] at (m.west) {\var{m}\textsubscript{1}};
  \foreach \i/\y in {1/1.4, 2/-1.4} {
    \node[q] (p\i) at (0,\y) {\var{P}};
    \draw[msg] (p\i) -- ($(q.west)+(0,{0.12*\y})$);
  }
  % un mensaje en camino
  \node[m] (mp) at (1.3,1.55) {};
  \draw[msg] (mp) -- (2.0,0.35);
  % el consumidor que tomó m₁ y falló
  \node[q] (cx) at (6.8,2.4) {\var{C}}; \tachar{cx}
  \node[m] (m2) at (4.95,2.45) {};
  \node[etiq, anchor=east] at (m2.west) {\var{m}\textsubscript{1}};
  \draw[msg] (cx.west) to[bend right=15] (4.6,0.6);
  % los otros dos, tirando de la cola
  \node[q] (c1) at (6.8,0.1) {\var{C}};
  \node[q] (c2) at (6.8,-2.1) {\var{C}};
  \draw[msg] (c1.west) -- ($(q.east)+(0.15,0.05)$);
  \draw[msg] (c2.west) to[bend left=12] ($(q.east)+(-0.1,-0.5)$);
  % el timer
  \draw[msg] (m.south) -- (3.4,-1.65);
  \begin{scope}[shift={(3.4,-2.05)}]
    \draw (0,0) circle (3.5mm);
    \draw[line width=0.9pt] (0,0) -- (0,0.24) (0,0) -- (0.16,-0.07);
  \end{scope}
  \node[align=center] at (3.4,-3.05) {visibility timer\\60 segundos};
  \node[anchor=north west, align=left] at (8.2,2.6) {– efecto balanceo\\\phantom{– }de carga};
\end{tikzpicture}\end{figura}
\end{document}
